\chapter{Gravitational Memory study}

\section{First meeting}

We talked in general about a couple of things. My next ``mini-project'' was to create a python code to read the strain off of the gravitational wave data and use that to create some audio. I have developed the code to work for hdf5 files and proved it worked with a black hole merger, I do not have the gravitational wave data for CCSN yet.

Also, the bigger project is to work on figuring out how to use 3Gen gravitational wave detectors to find CCSN. They even estimate these new projects may be able to see throughout the Virgo Cluster. We are presumably basing our work off of the paper \textbf{Detecting Gravitational Wave Memory in the Next Galactic Core-Collapse Supernova}

\subsection{First go: Detecting Gravitational Wave Memory in the Next Galactic Core-Collapse Supernova}

General notes from my first look at the paper detecting Gravitational Wave Memory in the Next Galactic Core-Collapse Supernova. Using data from GWOSC and several CCSN simulations they use some kind of linear prediction filtering system to develop what the waves should look like and how far away we could see them.

\subsubsection{Action Notes}

\begin{itemize}
  \item Learn Linear Prediction Filtration
  \item Learn librosa python module
  \item What is a whitened signal
  \item What is amplitude spectral density
\end{itemize}

I found a MITOpenCourseWare on signal theory that I will likely try to do

\subsubsection{Code}

Got some audio stuff working

\subsubsection{General Readings/Videos}

Convolution is defined via
\[
\sum_{k=-\infty}^\infty x[k]h[x-k] = (x \star h)[n]
  \]

  There is also a lot of other information listed but I am not sure what is important, the equations are generally pretty simple so I think I will just look them up.
\section{September 10}

add whiten to gw, zero out wavelets

Looking into trying without interpolation.

Look into generating gravitational waves from a 2d checkpoint file, calculate quadropole moments. Density variable multiply by coordinates. Stencile for time derivative or look into Pedro.

\subsection{Code}

I got the code for the audio of the CCSN. I did successfully whiten the gravitational wave data but was unable to get the black hole merger right. I verified what I did using the graph though

\subsection{Quadrupoles}

Calculating the quadrupole mass moment is fairly simple
\[
I_{2m} = \frac{16 \sqrt{3} \pi}{15} \int \rho Y_{2m}^{\star} r^2 dV
  \]

  The star superscript means we are using the complex conjugate.


  Rather than write everything out, calculate mass moment very simply, find its second order time derivative, then plug that into the gravitational wave equation.
